% ! This solution does not require any extra packages like etoolbox or xstring.

\begin{table*}[hbtp]
  \caption{Performance comparison of our proposed method with other fire/smoke detection methods on D-Fire \citeproc{ref-dfiredataset}{{[}xx{]}} and our self-collected dataset.}
  \label{tb:perf_results}
  \csvreader[separator=semicolon,
    no head,
    range=2-,
    centered tabularray = {
        width = 0.9\textwidth,
        colspec = {ccccccccccc},
        colsep = {5pt},
        hlines,
        vlines,
        cell{1}{1} = {r=2}{halign=c},
        cell{1}{2} = {c=5}{halign=c},
        cell{1}{7} = {c=5}{halign=c},
        cells = {font = \footnotesize, halign = c, valign = m},
        row{1} = {font = \bfseries\footnotesize},
        row{2} = {font = \bfseries\itshape\footnotesize},
    },
  ]{4.table/tb_perf_results.csv}{}{
    % Method (col1) + citation (col3)
    \ifcsvfirstrow
      {\csvexpval\csvcolii} % Header row: just the name
      { % Data rows: use \ifx for the most reliable empty check
        \ifx\csvcolv\empty
          {\csvexpval\csvcolii} % If \csvcolv is identical to \empty, print only the name
        \else
          {\csvexpval\csvcolii~\csvexpnot\citeproc{\csvexpval\csvcolv}{{[}\csvexpval\csvcolvi{]}}} % Otherwise, print the name + citation
        \fi
      } &
    % Other columns
    \csvexpval\csvcolvii &
    \csvexpval\csvcolviii &
    \csvexpval\csvcolix &
    \csvexpval\csvcolx &
    \csvexpval\csvcolxi &
    \csvexpval\csvcolxii &
    \csvexpval\csvcolxiii &
    \csvexpval\csvcolxiv &
    \csvexpval\csvcolxv &
    \csvexpval\csvcolxvi
  }
\end{table*}